% MIT OpenCourseWare: https://ocw.mit.edu
% 18.783 / 18.7831 Elliptic Curves Spring 2021
% License: Creative Commons BY-NC-SA 
% For information about citing these materials or our Terms of Use, visit: https://ocw.mit.edu/terms.
\newif\ifOCW
\OCWtrue
\documentclass[11pt]{article}
\usepackage{amsmath,amssymb,amsthm}
\usepackage{hyperref}
\hypersetup{colorlinks=true,urlcolor=blue,citecolor=blue,linkcolor=blue}
\ifOCW\usepackage{soul}\let\oldhref\href\renewcommand{\href}[2]{\oldhref{#1}{\ul{#2}}}\fi
\ifOCW\newcommand{\due}[1]{\hfill\phantom{Due: #1}}\else\newcommand{\due}[1]{\hfill{Due: #1}}\fi
\usepackage{courier}
\usepackage{graphicx}
\usepackage{array}
\usepackage{color}
\usepackage{enumerate}
\usepackage{nicefrac}
\usepackage{listings}
\lstset{
	basicstyle=\small\ttfamily,
	keywordstyle=\color{blue},
	language=python,
	xleftmargin=16pt,
}

\textwidth=5.8in
\textheight=9in
\topmargin=-0.5in
\headheight=0in
\headsep=.5in
\hoffset  -.4in
\pagestyle{plain}


% growing list of useful macros, use these where appropriate and add to this list as needed
\newcommand{\kbar}{\bar{k}}
\newcommand{\Fp}{\mathbb{F}_p}
\newcommand{\Fpbar}{\overline{\mathbb{F}}_p}
\newcommand{\Fq}{\mathbb{F}_q}
\newcommand{\Fqn}{\mathbb{F}_{q^n}}
\newcommand{\Fqbar}{\overline{\mathbb{F}}_q}
\newcommand{\F}{\mathbb{F}}
\newcommand{\Q}{\mathbb{Q}}
\newcommand{\Qbar}{\overline{\mathbb{Q}}}
\newcommand{\R}{\mathbb{R}}
\newcommand{\C}{\mathbb{C}}
\renewcommand{\H}{\mathbb{H}}
\newcommand{\D}{\mathbb{D}}
\renewcommand{\P}{\mathbb{P}}
\newcommand{\Z}{\mathbb{Z}}
\newcommand{\ZnZ}{\Z/n\Z}
\newcommand{\M}{\textsf{M}}
\newcommand{\Aut}{{\rm Aut}}
\newcommand{\Gal}{{\rm Gal}}
\newcommand{\GL}{{\rm GL}}
\newcommand{\PGL}{{\rm PGL}}
\newcommand{\End}{{\rm End}}
\newcommand{\DFT}{{\rm DFT}}
\newcommand{\dy}{\,dy}
\newcommand{\dx}{\,dx}
\newcommand{\tr}{\operatorname{tr}}
\newcommand{\kron}[2]{\bigl(\frac{#1}{#2}\bigr)}
\newcommand{\lcm}{\operatorname{lcm}}
\newcommand{\softO}{\widetilde{O}}
\newcommand{\ceil}[1]{\lceil{#1}\frac{1}eil}
\newcommand{\Exp}{{\rm E}}
\newcommand{\K}{\mathcal{K}}
\renewcommand{\O}{\mathcal{O}}
\newcommand{\OK}{\O_K}
\newcommand{\T}{{\rm T}}
\newcommand{\N}{{\rm N}}
\newcommand{\re}{\operatorname{re}}
\newcommand{\im}{\operatorname{im}}
\newcommand{\ord}{{\rm ord}}
\newcommand{\SL}{{\rm SL}}
\newcommand{\PSL}{{\rm PSL}}
\newcommand{\EllO}{{\rm Ell}_\O}
\newcommand{\fraka}{\mathfrak{a}}
\newcommand{\frakb}{\mathfrak{b}}
\newcommand{\frakc}{\mathfrak{c}}
\newcommand{\cl}{{\rm cl}}
\newcommand{\disc}{{\rm disc}}
\newcommand{\abcd}{\left(\begin{smallmatrix}a&b\\c&d\end{smallmatrix}\right)}
\newcommand{\ABCD}{\left(\begin{smallmatrix}A&B\\C&D\end{smallmatrix}\right)}
\newcommand{\p}{\mathfrak{p}}
\renewcommand{\a}{\mathfrak{a}}
\renewcommand{\b}{\mathfrak{b}}
\newcommand{\Frob}{{\rm Frob}}
\newcommand{\lt}{{\rm lt}}
\newcommand{\id}{{\rm id}}
\newcommand{\sumprime}{\sideset{}{'}\sum}

\newtheorem*{theorem}{Theorem}

\begin{document}
\setlength{\unitlength}{1in}
\begin{center}
\large
\textbf{18.783 Elliptic Curves\hspace{220pt}Spring~2021}\\\vspace{4pt}
\textbf{Problem Set \#8\due{04/21/2021}}\\\vspace{-6pt}
\normalsize
\begin{picture}(5.8,.1) 
\put(0,0) {\line(1,0){5.8}}
\end{picture}
\end{center}

\noindent\textbf{Description}: These problems are related to the material covered in Lectures 14--16.
\medskip

\noindent\textbf{Instructions}: Pick any combination of problems to solve that sums to 100 points.
Your solutions are to be written up in latex and submitted as a pdf-file\ifOCW\else to \href{https://www.gradescope.com/courses/238945}{Gradescope}\fi.

Collaboration is permitted/encouraged, but you must identify your collaborators or your group\ifOCW\else\ on \href{https://psetpartners.mit.edu}{pset partners}\fi, as well any references you consulted that are not listed in the \ifOCW\href{https://ocw.mit.edu/courses/mathematics/18-783-elliptic-curves-spring-2021/syllabus}{syllabus}\else\href{http://math.mit.edu/classes/18.783/2021/syllabus.html}{syllabus}\fi\ or \ifOCW\href{https://ocw.mit.edu/courses/mathematics/18-783-elliptic-curves-spring-2021/lecture-notes-and-worksheets}{lecture notes}\else\href{http://math.mit.edu/classes/18.783/2021/lectures.html}{lecture notes}\fi. If there are none write ``\textbf{Sources consulted: none}'' at the top of your solutions.  Note that each student is expected to write their own solutions; it is fine to discuss the problems with others, but your writing must be your own.

The first person to spot each non-trivial typo/error in the problem sets or lecture notes will receive 1-5 points of extra credit.

In cases where your solution involves writing code, please either include your code in your write up (as part of the pdf), or the name of a notebook in your 18.783 CoCalc project containing you code (please use a separate notebook for each problem).


\subsection*{Problem 1. From lattices to elliptic curves (49 points)}
In this problem you will explicitly construct an elliptic curve corresponding to a given lattice $L=[1,\tau]$ by computing the $j$-invariant
\[
j(L) = 1728\frac{g_2(L)^3}{g_2(L)^3-27g_3(L)^2},
\]
where
\begin{equation}\label{eq:g2g3}
g_2(L)=60\sum_{\substack{m,n\in\Z\\(m,n)\ne (0,0)}}\frac{1}{(m+n\tau)^4},\qquad\text{and}\qquad g_3(L)=140\sum_{\substack{m,n\in\Z\\(m,n)\ne(0,0)}}\frac{1}{(m+n\tau)^6}.
\end{equation}

Let $\tau=(1+\sqrt{-7})/2$, so that the lattice $L=[1,\tau]$ is the ring of integers of $\Q(\sqrt{-7})$.

\begin{enumerate}[{\bf (a)}]
\item 
As a warmup, use Sage to approximate the complex numbers $g_2(L)$ and $g_3(L)$ by summing over lattice points $m+n\tau$ with $|m|,|n|< 10$.
To work over the complex numbers, you will need to set \texttt{tau=CC((1+sqrt(-7))/2)} (by default, Sage will work with a symbolic representation of $\tau$, the wrapper \texttt{CC( )} coerces $\tau$ to $\C$).
Now define the elliptic curve $E:y^2=x^3-g_2(L)/4 x - g_3/4(L)$ over $\C$ using your approximations of $g_2(L)$ and $g_3(L)$ (use \texttt{E=EllipticCurve([-g2/4,-g3/4])}).
Let $z=0.N$, where $N$ is the last 4 digits of your student ID, and compute $x=\wp(z;L)$ and $y=\wp'(z;L)$ by summing over lattice points $m+n\tau$ with $|m|,|n|<10$.
The point $(x,y/2)$ will then be approximately on the elliptic curve $E$ that you defined.
To get an exact point $P\in E(\C)$, use \texttt{E.lift\_x(x)}, which will cause sage to choose an exact $y$-value corresponding to $x$ (with an arbitrary sign choice).
You should find that $y$-coordinate of $P$ is approximately $\pm y/2$.
Now compute $7P$ in Sage, and compare its $x$ and $y$ coordinates with $\wp(7z)$ and $\wp'(7z)/2$ (the sign on the $y$-coordinate may be off, don't worry about this); this corresponds to comparing $7\Phi(z)$ with $\Phi(7z)$, where $\Phi$ is the isomorphism $\C/L\to E(\C)$ defined by $z\mapsto (\wp(z),\wp'(z)/2)$.
In your answer report the values of $g_2(L)$ and $g_3(L)$ that you computed, along with $z$, $P$, $(\wp(z),\wp'(z)/2)$, $7P$, and $(\wp(7z),\wp'(7z)/2)$.

\item
While it will not be apparent from your work in part (a), the elliptic curve $E$ corresponding to $L=[1,\tau]$ is actually defined over $\Q$, even though $g_2(L)$ and $g_3(L)$ are not.
To see this we need to compute the $j$-invariant
\[
j(L) = 1728\frac{g_2(L)^3}{g_2(L)^3-27g_3(L)^2}.
\]
Use Sage to approximate $j(L)$ by computing the sums for $g_2(L)$ and $g_2(L)$ over lattice points with $|m|,|n|< r$ for increasing values of $r=10,20,30,\ldots$, until you are convinced you can correctly approximate the real and imaginary parts of $j(L)$ to one decimal place. You should find that $j(L)$ approaches an integer as $r$ increases.
In your solution list the value of $r$ you used, the complex value of $j(L)$ you computed (to 8 decimal places), as well as the integer obtained.

\item
Use the $j$-invariant $j=j(L)$ you computed in part (b) to construct an elliptic curve $E\colon y^2 = x^3+Ax+B$ over $\Q$ using
$A=3j(1728-j)$ and $B=2j(1728-j)^2$.
Heuristically verify that $\End^0(E)\simeq \Q(\sqrt{-7})$ by checking for each prime $p<1000$ where $E$ has good reduction that:
\begin{enumerate}
\item[(i)] if $\bigl(\frac{-7}{p}\bigr)=1$ then $4p=t^2+7v^2$ for some $v\in \Z$, where $t=\tr\pi_{E_p}$,
\item[(ii)] if $\bigl(\frac{-7}{p}\bigr)=-1$ then $t=\tr\pi_{E_p}\equiv 0\bmod p$, so $E_p$ is supersingular,
\end{enumerate}
where $E_p/\F_p$ is the elliptic curve given by reducing the equation for $E$ modulo $p$.

\end{enumerate}

\noindent
Using \eqref{eq:g2g3} to approximate $g_2(L)$ and $g_3(L)$ is inefficient because the sums converge very slowly;
a better approach is to use their $q$-expansions.  If we put $q=\exp(2\pi i \tau)$ then
\begin{align*}
g_2([1,\tau])&=\frac{4\pi^4}{3}\left(1 + 240 \sum_{k=1}^\infty\frac{k^3q^k}{1-q^k}\right)\quad\text{and}\quad
g_3([1,\tau])&=\frac{8\pi^6}{27}\left(1 - 504 \sum_{k=1}^\infty\frac{k^5q^k}{1-q^k}\right);
\end{align*}
see \cite[p.\ 275]{washington}.

\begin{enumerate}[{\bf (a)}]
\setcounter{enumi}{3}
\item
Repeat part (b) using the $q$-expansion formulas for $g_2(L)$ and $g_3(L)$, truncating the sums after 1000 terms.
Extend the precision of your computations by defining \texttt{CC=ComplexField(500)}, and use
\texttt{q=CC(exp(2*pi*sqrt(-1)*tau))} to compute $q$  (\textbf{important}: use \texttt{tau=(1+sqrt(-7))/2}, coercing $\tau$ to $\C$ before computing $q$ will result in a loss of precision).
Compare the resulting approximation to $j(L)$ to the one you computed in part (b) by listing the the real and imaginary parts of both approximations to $8$ decimal places.

\item
Use your improved algorithm to compute the $j$-invariant of the lattice $L=[1,\sqrt{-7}]$.
Assuming it is a rational integer, construct the corresponding elliptic curve and heuristically verify that it also has CM by  $\Q(\sqrt{-7})$ (note that in this case $L$ is not the maximal order in $\Q(\sqrt{-7})$).

\item
Now let $L=[1,(1+\sqrt{-23})/2]$ be the ring of integers of $\Q(\sqrt{-23})$.
After approximating $j(L)$ you will find that it does not appear to be a rational integer.
But it is an algebraic integer.  Use Sage to find its minimal polynomial using the \texttt{algdep} method with a degree bound of 4 and the optional \texttt{use\_digits} parameter set to 100  (you should get a monic polynomial of degree 3; if not, you have made a mistake or are not using enough precision).

\item
Let $H(x)$ be the minimal polynomial you computed in part (f) and let $D=-23$.
Find a prime $p$ for which $\bigl(\frac{D}{p}\bigr)=1$ and $H(x)$ splits completely into linear factors in $\Fp[x]$, and let $r$ be one of its roots.
Construct an elliptic curve $E/\Fp$ with $j$-invariant~$r$ and compute its trace of Frobenius~$t$.
Verify that $4p=t^2-v^2D$ for some integer $v$.
Repeat this verification for every prime $p<1000$ for which $\bigl(\frac{D}{p}\bigr)=1$ and $H(x)$ splits completely in $\Fp[x]$.
Now use this method to construct an elliptic curve with CM by $\Q(\sqrt{D})$ over a 256-bit finite field.

\end{enumerate}

\subsection*{Problem  2. From elliptic curves to lattices (49 points)}
We now consider the problem of determining the lattice $L$ corresponding to an elliptic curve $E\colon y^2=x^3+Ax+B$. This is known as ``computing the periods" of $E$, and involves computing approximate solutions to certain elliptic integrals associated to $E$, as explained in \cite[\S 9.4]{washington}.
To simplify matters, we will focus on the case where $A$ and $B$ are real numbers.

Given two positive real numbers $a$ and $b$, define the sequences $\{a_n\}$ and $\{b_n\}$ as follows:
\begin{equation}\label{eq:agm}
a_0 = a,\qquad b_0 = b,\qquad
a_n = \frac{a_{n-1}+b_{n-1}}{2},\qquad
b_n = \sqrt{a_{n-1}b_{n-1}}.
\end{equation}
As proven in \cite[Prop.\ 9.23]{washington}, these sequences both converge to a common limit $M(a,b)$, which is defined as the \emph{arithmetic-geometric mean} (AGM) of $a$ and $b$.
As with Newton iteration, the rate of convergence is doubly exponential, which makes the arithmetic-geometric mean a powerful tool for numerical algorithms.

When the cubic $f(x)=x^3+Ax+B$ has three real roots $e_1<e_2 < e_3$, we can compute a lattice $L=[\omega_1,\omega_2]$ for $E$ via the formulas
\begin{align*}
\omega_1 &= \frac{\pi }{M(\sqrt{e_3-e_1},\sqrt{e_3-e_2})},\\
\omega_2 &= \frac{\pi i}{M(\sqrt{e_3-e_1},\sqrt{e_2-e_1})},
\end{align*}
as proven in \cite[Thm.\ 9.26]{washington}.
When $f(x)=x^3+Ax+B$ has just one real root $e_1$, we let $e_2=\sqrt{3e_1^2+A}$ and use the formulas
\begin{align*}
\omega_1 &= \frac{2\pi }{M(2\sqrt{e_2},\sqrt{2e_2+3e_1})},\\
\omega_2 &=-\frac{\omega_1}{2}+ \frac{\pi i}{M(2\sqrt{e_2},\sqrt{2e_2-3e_1})}.
\end{align*}
The resulting lattice $L=[\omega_1,\omega_2]$ then satisfies $g_2(L)=-4A$ and $g_3(L)=-4B$, so that the elliptic curve $y^2=4x^3-g_2(L)x-g_3(L)$ corresponding to the torus $\C/L$ is isomorphic to our original curve $E$.


\begin{enumerate}[{\bf (a)}]
\item Implement an algorithm in Sage to approximate $M(a,b)$ using \eqref{eq:agm}.
Using your algorithm, compute the RHS of the identity\footnote{This identity was of great interest to Gauss; the quantity $1/M(1,\sqrt{2})=0.8346268\ldots$ is known as Gauss's constant.  A proof can be found in \cite[Ex.\ VI.6.12-14]{silverman} (NB: there is a typo in part (f) of Exercise VI.6.14 in \cite{silverman}: the quantity $M(1,\sqrt{2})$ should appear in the denominator, as above).}
\[
\int_0^1 \frac{dz}{\sqrt{1-z^4}} = \frac{\pi}{2M(1,\sqrt{2})},
\]
use Sage to compute the LHS, and verify that these values agree to, say, 100 decimal places.
You will need to extend the precision of the real field to do this: use \texttt{RR=RealField(1000)} to get 1000 bits of precision, and then be sure to coerce the arguments to $M(a,b)$ into \texttt{RR} using \texttt{M(RR(a),RR(b))}.
To compute the LHS in Sage, use \texttt{RR(integral(1/sqrt(1-x**4),x,0,1))}.

\item
Using the formulas above, approximate the periods $\omega_1$ and $\omega_2$ associated to the elliptic curve $E\colon y^2=x^3-35x-98$.
Compute the ratio $\tau = \omega_2/\omega_1$, so that $L$ is homothetic to $[1,\tau]$, and then compute the $j$-invariant $j(L)$ and compare it to $j(E)$.
Attempt to identify $\tau$ as an algebraic number in a quadratic field using the \texttt{algdep} method in Sage, with the degree bound set to 2.
In your answers, just list your values for $\omega_1,\omega_2,\tau$, and $j(L)$ out to 16 decimal places, even though you may need to use higher precision in your computations, and list the polynomial computed by \texttt{algdep}.

\item Do the same thing for the elliptic curves
\begin{align*}
E_1\colon y^2&=x^3-7x+6,\\
E_2\colon y^2&=x^3-608x+5776,\\
E_3\colon y^2 &= x^3 - 34790720x + 78984748304.
\end{align*}
In cases where you are able to provisionally identify $\tau$ as an algebraic number in a quadratic field $K=\Q(\sqrt{D})$, heuristically test whether $E_i$ has CM by $K$ by checking if it has supersingular reduction modulo good primes $p<1000$ for which $\bigl(\frac{D}{p}\bigr)=-1$.
For the $E_i$ where this test is successful, it may be that $\tau$ is an algebraic number but not an algebraic integer.
Show that in each such case $\tau$ is equivalent under the action of $\SL_2(\Z)$ to an algebraic integer with real part $0$ or $-1/2$.

\item Lastly, compute the periods for an elliptic curve that is defined over $\R$ but not over~$\Q$ (or any number field).
Let $N$ be the last 4 digits of your student ID, and compute the periods for $E\colon y^2 = x^3 + \pi x + N$,
where $\pi=3.1415\ldots$ is transcendental.
\end{enumerate}

\subsection*{Problem  3.  Elliptic curves over $\boldsymbol{\R}$ (49 points)}
Let $L=[1,\tau]$ be a lattice with $0\le \re \tau < 1$ (note that every lattice in $\C$ is homothetic to such a lattice $L$), let $E/\C$ be the corresponding elliptic curve $$y^2=4x^3-g_2(L)x-g_3(L),$$
and let $\Phi(z)=(\wp(z),\wp'(z))$ be the isomorphism from $\C/L$ to $E(\C)$.

\begin{enumerate}[{\bf (a)}]
\item Prove that $E$ is defined over $\R$ (meaning $g_2(L),g_3(L)\in\R$) if and only if $L$ is stable under complex conjugation.

\item Characterize the lattices $L=[1,\tau]$ that are stable under complex conjugation by giving necessary and sufficient conditions on $\tau$ (with $0\le \re \tau < 1$ as above).

\item Does your answer to (b) cover every lattice $L=[1,\tau]$ with $j(L)\in \R$?
If not, describe those that are missing and why these do not contradict your answers to (a) and (b).
\end{enumerate}

Let us now assume that $E$ is defined over $\R$ (so $g_2(L),g_3(L)\in\R$), but still view $E$ as an elliptic curve over~$\C$, using $E(\R)$ to denote the subgroup of real points in $E(\C)$.

\begin{enumerate}[{\bf (a)}]
\setcounter{enumi}{3}
\item Prove that if $z\in\R$ then $\Phi(z)\in E(\R)$.
\item Prove that if $z=\frac{1}{2}+it$ with $t\in\R$ then $\wp(z)\in \R$ but $\wp'(z)\not\in \R$, unless $\wp'(z)=0$.
Conclude that there is at most one $z_0\not\in \R$ in the  fundamental parallelogram $\mathcal{F}_0$ for $L$ with real part $1/2$ for which $\Phi(z_0)\in E(\R)$.  Show that such a $z_0$ exists if and only if the cubic $f(x)=4x^3-g_2(L)x-g_3(L)$ has three real roots.
\item Prove that the pre-image $\Phi^{-1}(E(\R))$ is given by
\[
\mathcal{R}_0 =
\begin{cases}
\{z\in\mathcal{F}_0: \im z = 0\}&\text{if $f(x)$ has one real root},\\
\{z\in\mathcal{F}_0:\im z = 0 \text{ or }\im z= \frac{1}{2}\im \tau\}&\text{if $f(x)$ has three real roots}.
\end{cases}
\]
Prove that $L$ is rectangular ($\re \tau = 0$) if and only  if we are in the latter case.

\item Conclude that $E(\R)$ has either one or two components, depending on whether its 2-torsion subgroup is cyclic or not.
Show that $E(\R)\simeq \R/\Z$ in the first case and $E(\R)\simeq \R/\Z\oplus \Z/2$ in the second.
\end{enumerate}

\subsection*{Problem 4.  The modular group $\boldsymbol{\SL_2(\Z)}$ (49 points)}
Let $\Gamma=\SL_2(\Z)$ and let $\H=\{z\in\C:\im z > 0\}$ be the upper half plane.
For each $\gamma=\abcd\in\Gamma$ and $\tau\in\H$, define
\[
\gamma\tau = \frac{a\tau+b}{c\tau+d}.
\]
Let $S=\left(\begin{smallmatrix}0&-1\\1&0\end{smallmatrix}\right)$ and let $T=\left(\begin{smallmatrix}1&1\\0&1\end{smallmatrix}\right)$.
\begin{enumerate}[{\bf (a)}]
\item Prove that $\Gamma$ is generated by $S$ and $T$.
\item Prove that $\gamma\tau\in\H$ for all $\gamma\in\Gamma$ and $\tau\in\H$.
\item Prove that the map from $\Gamma\times\H$ to $\H$ that sends $(\gamma,\tau)$  to $\gamma\tau$ is a group action.
\item Compute the stabilizers of $i:=e^{\pi i/2}$ and $\rho:=e^{2\pi i/3}$ under the action of $\Gamma$.
Express the elements of each stabilizer in terms of $S$ and $T$.
\item Prove that  the stabilizer of every element of $\H$ that is not $\Gamma$-equivalent to $i$ or $\rho$ is the subgroup of order 2 consisting of $\pm I$, where $I$ is the $2\times 2$ identity matrix.
\end{enumerate}

The \emph{extended upper half plane} $\H^*$ is defined as $\H\cup \P^1(\Q)$, where $\P^1(\Q)$ is the projective line over $\Q$,
consisting of all projective points $(x:y)$ with integer coordinates.
One can view $\P^1(\Q)$ as $\Q\cup\{\infty\}$, where $\Q$ consists of the points $x/y=(x:y)$ with $y\ne 0$ and $\infty$ is the point $(1:0)$ ``at infinity".
We extend the action of $\Gamma$ to $\H^*$ by defining
\[
\gamma(x:y)=(ax+by:cx+dy)
\]
for each $\gamma=\abcd\in \Gamma$ and $(x:y)\in\P^1(\Q)$.

\begin{enumerate}[{\bf(a)}]
\setcounter{enumi}{5}
\item Prove that the elements of $\P^1(\Q)$ are all $\Gamma$-equivalent.
\item Compute the stabilizers of $0=(0:1)$ and $\infty=(1:0)$.
\item Compute the stabilizer of $(x:y)$ when $xy\ne 0$.
\end{enumerate}

\subsection*{Problem 5. Survey (2 points)}

Complete the following survey by rating each of the problems you attempted on a scale of 1 to~10 according to how interesting you found it (1 = ``mind-numbing," 10 = ``mind-blowing"), and how difficult you it (1 = ``trivial," 10 = ``brutal").  Also estimate the amount of time you spent on each problem to the nearest half hour.

\begin{center}
\begin{tabular}{l|r|r|r|}
& Interest & Difficulty & Time Spent\\\hline
Problem 1 & & & \\\hline
Problem 2 & & & \\\hline
Problem 3 & & & \\\hline
Problem 4 & & & \\\hline
\end{tabular}
\end{center}
\noindent
Also, please rate each of the following lectures that you attended, according to the quality of the material (1=``useless", 10=``fascinating"), the quality of the presentation (1=``epic fail", 10=``perfection"), the pace (1=``way too slow", 10=``way too fast", 5=``just right") and the novelty of the material (1=``old hat", 10=``all new").

\begin{center}
\begin{tabular}{l|l|r|r|r|r|r}
Date & Lecture Topic & Material & Presentation & Pace & Novelty\\\hline
4/12 & Elliptic curves over $\C$ (part 2) & & & & \\\hline 
4/14 & Complex multiplication (CM) & & & & \\\hline 
\end{tabular}
\end{center}

\noindent
Please feel free to record any additional comments you have on the problem sets or lectures, in particular, ways in which they might be improved.

\begin{thebibliography}{9}
\bibitem{silverman} J.H. Silverman, \href{https://link.springer.com/book/10.1007/978-0-387-09494-6}{\textit{The arithmetic of elliptic curves}}, second edition, Springer, 2009;.
\bibitem{washington} L.C. Washington, \href{https://doi.org/10.1201/9781420071474}{\textit{Elliptic curves: Number theory and cryptography}}, second edition, Chapman and Hall/CRC, 2008. 
\end{thebibliography}
\end{document}